\Article{15}
Ist $X$ durch eine gegebene Funktion von $u$ und $\lambda$ ausgedrückt, so lässt sich daraus $Y$ a priori ableiten. Man setze das Integral $\int_0^u X\mathrm{d}u=T$ und halte $\lambda$ bei der Integration konstant. Differenziert man dann auf dieselbe Weise nach $u$, so gilt offenbar $\frac{\mathrm{d}(V+RT)}{\mathrm{d}u}=0$. Damit ist $V+RT$ eine von $u$ unabhängige Größe. Mit anderen Worten ist sie an allen Punkten eines einzelnen Meridians konstant. Sie muss daher auch überall konstant sein, weil alle Meridiane in den Polen zusammenlaufen. Setzt man den Wert von $V$ am Nordpol $=V^*$, so ergibt sich also
\[
T=\frac{V^*-V}{R}
\]
und damit
\[
Y=\frac{\mathrm{d}T}{\sin u\cdot\mathrm{d}\lambda}
\]
Das Ergebnis lässt sich auch so ausdrücken:
\[
Y=\frac{1}{\sin u}\int_0^u\frac{\mathrm{d}X}{\mathrm{d}\lambda}\cdot\mathrm{d}u
\]

\Article{16}
Dieser bemerkenswerte Satz, dass \textit{wenn die nach Norden gerichtete Komponente der horizontalen magnetischen Kraft für die gesamte Erdoberfläche gegeben ist, die nach Westen (oder Osten) gerichtete Komponente von selbst daraus folgt}, gilt umgekehrt nur mit einer Abänderung. Ist nämlich $Y$ durch eine gegebene Funktion von $u$ und $\lambda$ ausgedrückt und bezeichnet man mit $U$ das unbestimmte Integral $\int\sin u\cdot Y\mathrm{d}\lambda$, wobei $u$ bei der Integration als konstant gilt, so ist $\frac{\mathrm{d}(V+RU)}{\mathrm{d}\lambda}=0$. oder $V+RU$ ist eine von $\lambda$ unabhängige Größe, mithin im Allgemeinen eine Funktion von $u$. Also ist auch $\frac{\mathrm{d}(V+RU)}{R\mathrm{d}u}=\frac{\mathrm{d}U}{\mathrm{d}u}-X$ eine solche Funktion. Die Formel $\frac{\mathrm{d}U}{\mathrm{d}u}$ liefert also einen unvollständigen Ausdruck für $X$, weil ein Bestandteil, der nur $u$ enthält, unbestimmt bleibt. Der fehlende Bestandteil lässt sich aber ergänzen, wenn außer dem Ausdruck für $Y$ auch der für $X$ auf irgendeinem bestimmten Meridian oder allgemeiner auf irgendeiner vom \SourceLeaf{144}Nordpol zum Südpol reichenden Linie bekannt ist. Man sieht also, dass \textit{wenn man die nach Westen gerichtete Komponente der horizontalen magnetischen Kraft für die gesamte Erdoberfläche und die nach Norden gerichtete Komponente für alle Punkte irgendeiner vom Nordpol zum Südpol verlaufenden Linie kennt, die letztere Komponente für die gesamte Erdoberfläche von selbst daraus folgt.}

\Article{17}
Die vorangegangenen Untersuchungen betreffen allein den horizontalen Teil der erdmagnetischen Kraft. Um auch den vertikalen Teil zu erfassen, müssen wir die Aufgabe in ihrer ganzen Allgemeinheit betrachten, also $V$ als Funktion von \textit{drei} veränderlichen Größen auffassen, die die Lage eines beliebigen Punkts $O$ im Raum angeben. Wir wählen dazu die Entfernung $r$ vom Erdmittelpunkt, den Winkel $u$, den $r$ mit dem nördlichen Teil der Erdachse bildet, und den Winkel $\lambda$ zwischen der durch $r$ und die Erdachse gelegten Ebene und einem festen Meridian, nach Osten als positiv gezählt.

Die Funktion $V$ sei in eine Reihe nach fallenden Potenzen von $r$ entwickelt, der wir die folgende Form geben
\[
V=\frac{RRP^0}{r}+\frac{R^3P'}{rr}+\frac{R^4P''}{r^3}+\frac{R^5P'''}{r^4}+\text{u. s. w.}
\]
Die Koeffizienten $P^0,P',P''$ usw. sind hier Funktionen von $u$ und $\lambda$. Um zu sehen, wie sie mit der Verteilung der magnetischen Flüssigkeit im Erdinnern zusammenhängen, sei $\mathrm{d}\mu$ ein Element dieser Flüssigkeit und $\rho$ seine Entfernung von $O$. Für $\mathrm{d}\mu$ sollen $r^0,u^0,\lambda^0$ dasselbe bedeuten wie $r,u,\lambda$ für $O$. Damit gilt $V=-\int\frac{\mathrm{d}\mu}{\rho}$, wobei sich die Integration über alle $\mathrm{d}\mu$ erstreckt. Ferner gilt
\[
\rho=\BareRoot{\{rr-2rr^0(\cos u\cos u^0+\sin u\sin u^0\cos(\lambda-\lambda^0))+r^0r^0\}}
\]
und entwickelt man $\frac1\rho$ in die Reihe
\[
\frac1\rho=\frac1r(T^0+T'\cdot\frac{r^0}{r}+T''\cdot\frac{r^0r^0}{rr}+\text{u. s. w.})
\]
so erhält man
\[
\begin{aligned}RRP^0&=-\int T^0\mathrm{d}\mu,\quad R^3P'=-\int T'r^0\mathrm{d}\mu,\\R^4P''&=-\int T''r^0r^0\mathrm{d}\mu\quad\text{u. s. w.}\end{aligned}
\]
Da $T^0=1$ ist, folgt aus der Grundannahme, dass in jedem messbaren Teilchen des Trägers und somit auch in der \textit{ganzen} Erde gleich viel positive wie negative Flüssigkeit vorhanden ist, oder dass $\int\mathrm{d}\mu=0$ gilt,
\SourceLeaf{145}
\[
P^0=0
\]
also fällt das erste Glied unserer Reihe für $V$ weg.

Man sieht ferner, dass $P'$ die Form
\[
R^3P'=\alpha\cos u+\beta\sin u\cos\lambda+\gamma\sin u\sin\lambda
\]
hat, wobei $\alpha=-\int r^0\cos u^0\cdot\mathrm{d}\mu$, $\beta=-\int r^0\sin u^0\cos\lambda^0\mathrm{d}\mu$ und $\gamma=-\int r^0\sin u^0\sin\lambda^0\cdot\mathrm{d}\mu$. Also sind $-\alpha,-\beta,-\gamma$ nach der in Artikel 5 der \textit{Intensitas vis magneticae} gegebenen Definition die Momente des Erdmagnetismus bezüglich dreier rechtwinklig zueinander stehender Achsen. Die erste ist die Erdachse, die zweite und die dritte sind die Äquatorradien für die Längen $0$ und $90^0$.

Die allgemeinen Formeln für alle Koeffizienten der Reihe für $\frac1\rho$ können wir als bekannt voraussetzen. Für unseren Zweck genügt die Bemerkung, dass die Koeffizienten in Bezug auf $u$ und $\lambda$ ganze rationale Funktionen von $\cos u$, $\sin u\cos\lambda$ und $\sin u\sin\lambda$ sind, und zwar $T''$ von zweiter Ordnung, $T'''$ von dritter Ordnung usw. Dasselbe gilt also auch für die Koeffizienten $P'',P'''$ usw.

Die Reihen für $\frac1\rho$ und für $V$ konvergieren, solange $r$ nicht kleiner als $R$ ist, oder genauer, nicht kleiner als der Radius einer Kugel, die sämtliche magnetischen Teile der Erde einschließt.

\Article{18}
Die Funktion $V$ genügt aufgrund ihrer Zusammensetzung aus $-\int\frac{\mathrm{d}\mu}{\rho}$ der folgenden partiellen Differentialgleichung
\[
0=\frac{r\mathrm{d}\mathrm{d}rV}{\mathrm{d}r^2}+\frac{\mathrm{d}\mathrm{d}V}{\mathrm{d}u^2}+\operatorname{cotg}u\cdot\frac{\mathrm{d}V}{\mathrm{d}u}+\frac1{\sin u^2}\cdot\frac{\mathrm{d}\mathrm{d}V}{\mathrm{d}\lambda^2}
\]
die eine Umformung der bekannten Gleichung
\[
0=\frac{\mathrm{d}\mathrm{d}V}{\mathrm{d}x^2}+\frac{\mathrm{d}\mathrm{d}V}{\mathrm{d}y^2}+\frac{\mathrm{d}\mathrm{d}V}{\mathrm{d}z^2}
\]
ist, wobei $x,y,z$ die rechtwinkligen Koordinaten von $O$ bezeichnen. Setzt man in die erstgenannte Gleichung den Wert von $V$
\[
V=\frac{R^3P'}{rr}+\frac{R^4P''}{r^3}+\frac{R^5P'''}{r^4}+\text{u. s. w.}
\]
ein, so zeigt sich, dass für die einzelnen Koeffizienten $P',P'',P'''$ usw. ebenfalls \SourceLeaf{146}partielle Differentialgleichungen gelten, deren allgemeiner Ausdruck lautet
\[
0=n(n+1)P^{\mathrm{n}}+\frac{\mathrm{d}\mathrm{d}P^{\mathrm{n}}}{\mathrm{d}u^2}+\operatorname{cotg}u\frac{\mathrm{d}P^{\mathrm{n}}}{\mathrm{d}u}+\frac1{\sin u^2}\cdot\frac{\mathrm{d}\mathrm{d}P^{\mathrm{n}}}{\mathrm{d}\lambda^2}
\]

Aus dieser Gleichung und der Bemerkung im vorangegangenen Artikel ergibt sich die allgemeine Form von $P^{\mathrm{n}}$. Bezeichnet man nämlich mit $P^{\mathrm{n,m}}$ die folgende Funktion von $u$
\[
\begin{aligned}
\{\cos u^{n-m}&-\frac{(n-m)(n-m-1)}{2(2n-1)}\cos u^{n-m-2}\\
&+\frac{(n-m)(n-m-1)(n-m-2)(n-m-3)}{2\cdot4(2n-1)(2n-3)}\cos u^{n-m-4}\\&-\text{u. s. w.}\}\sin u^m
\end{aligned}
\]
so hat $P^{\mathrm{n}}$ die Form einer Summe aus $2n+1$ Teilen
\[
\begin{aligned}
P^{\mathrm{n}}&=g^{\mathrm{n,0}}P^{\mathrm{n,0}}+(g^{\mathrm{n,1}}\cos\lambda+h^{\mathrm{n,1}}\sin\lambda)P^{\mathrm{n,1}}\\
&\quad+(g^{\mathrm{n,2}}\cos2\lambda+h^{\mathrm{n,2}}\sin2\lambda)P^{\mathrm{n,2}}+\text{etc.}+(g^{\mathrm{n,n}}\cos n\lambda+h^{\mathrm{n,n}}\sin n\lambda)P^{\mathrm{n,n}}
\end{aligned}
\]
wobei $g^{\mathrm{n,0}},g^{\mathrm{n,1}},h^{\mathrm{n,1}},g^{\mathrm{n,2}}$ usw. bestimmte Zahlenkoeffizienten sind.

